\section{综合讨论}

\subsection{三任务横向对比}

三项任务共享同一数据集、同一预处理流水线与同一套验证框架，但呈现出高度一致的结构特征：稀疏物理公式即可逼近预测上限。表\ref{tab:overall}汇总了关键指标。

\begin{table}[H]
\centering
\caption{三项任务综合对比}
\label{tab:overall}
\begin{tabular}{lcccl}
\toprule
\textbf{维度} & \textbf{任务一} & \textbf{任务二} & \textbf{任务三} & \textbf{说明} \\
\midrule
核心驱动变量 & 周里程 & 能耗 $\times$ 电价 & 日通勤 & \\
最终模型 & 比率均值 & 级联公式 & 过原点 OLS & \\
adj-$R^2$ & 0.876298 & 0.917817 & 0.905416 & 5折$\times$3种子均值 \\
折间标准差 & 0.000007 & 0.000004 & 0.000002 & 高度稳定 \\
残差树模型 $\Delta R^2$ & $-0.001430$ & $-0.000922$ & $-0.000355$ & 全部为负 \\
特征块消融 & 全部下降 & 级联最优 & 全部下降 & 稀疏最优 \\
\bottomrule
\end{tabular}
\end{table}

三项任务的共同结论是：数据由近似确定性的物理关系生成，稀疏模型已达预测上限，复杂模型（浅层提升树）在 Bootstrap 置信区间下全面劣化。

\subsection{物理结构优先的方法论}

本研究与初赛阶段形成鲜明对照。初赛三任务均为分类，采用 CatBoost 统一建模并取得了较强的信号；而本次复赛的回归目标背后存在近似确定性的生成公式，此时任何数据驱动模型都难以超越对该公式的精确刻画。这一对照指向一个可推广的方法论原则：\textbf{在投入大规模模型搜索之前，先通过比率诊断与相关分析识别数据的底层结构}，结构存在时以公式/线性模型为主，仅在残差存在稳定信号时才引入非线性模型。

\subsection{调整决定系数下的稀疏性权衡}

复赛以 adj-$R^2$ 为评分指标，这与初赛的准确率存在本质差异：adj-$R^2$ 会对新增预测变量施加自由度惩罚。因此，一个无增益特征的加入，即使不改变 $R^2$，也会降低 adj-$R^2$。本文的消融实验反复印证了这一点：特征块越多，adj-$R^2$ 越低。在``特征越多越好''的惯性思维下，这一指标反而奖励了最简洁的模型，与本文发现的数据结构不谋而合。

\subsection{严格级联的防泄漏价值}

任务二的严格嵌套 OOF 级联是本文在方法论层面的核心贡献之一。若图省事直接在训练时使用真实能耗、测试时改用预测能耗，将造成 teacher forcing，线下指标虚高、线上无法复现。本文的嵌套方案从机制上保证了能耗预测与成本建模的信息隔离，可复用于任何``预测-再预测''的级联回归场景。

\subsection{负结果的方法论价值}

与初赛任务二类似，本文保留了完整的负结果记录：题面建议的车型/地域交互、浅层提升树模型均被证伪。这些负结果的意义在于：

\begin{enumerate}[itemsep=0pt]
\item 明确界定了``模型选择不当''与``数据结构本身简单''两种情况的边界；
\item 通过 Bootstrap 置信区间而非单次实验的偶然提升来判定模型取舍，避免了过度拟合噪声；
\item 为同类合成数据竞赛提供了一套``结构识别 $\to$ 稀疏基线 $\to$ 消融 $\to$ 残差诊断 $\to$ 最终确认''的可复用流程。
\end{enumerate}
